\section*{PHẦN MỞ ĐẦU}
\addcontentsline{toc}{section}{PHẦN MỞ ĐẦU}

\subsection*{1. Tính cấp thiết của đề tài nghiên cứu}

Giáo dục nghề nghiệp giữ vị trí then chốt trong cấu trúc hệ thống giáo dục quốc dân của mọi quốc gia, đóng vai trò trực tiếp trong việc chuyển hóa tiềm năng dân số thành nguồn nhân lực kỹ thuật phục vụ trực tiếp cho quá trình sản xuất, lưu thông vật chất và cung ứng dịch vụ của nền kinh tế. Đối với Việt Nam, xuyên suốt hơn tám thập kỷ kể từ ngày lập quốc mùa thu năm 1945 đến nay, lịch sử phát triển của hệ thống đào tạo nghề và giáo dục nghề nghiệp luôn gắn bó máu thịt với vận mệnh của dân tộc, phản ánh sinh động các giai đoạn chuyển mình từ chiến tranh vệ quốc gian khổ sang thời kỳ kiến thiết kinh tế xã hội chủ nghĩa, từ cơ chế kế hoạch hóa tập trung bao cấp sang kinh tế thị trường định hướng xã hội chủ nghĩa và hội nhập quốc tế sâu rộng.

Tại mỗi khúc quanh của lịch sử, tổ chức bộ máy quản lý nhà nước, mục tiêu đào tạo, cơ cấu trình độ và phương thức vận hành của mạng lưới cơ sở đào tạo nghề đều có những bước biến chuyển sâu sắc nhằm thích ứng với đòi hỏi của hiện thực khách quan. Tuy nhiên, tiến trình phát triển này không diễn ra theo một đường thẳng thuần túy mà đan xen nhiều bước thăng trầm, xuất hiện không ít cuộc chuyển giao thể chế phức tạp, tồn tại những xung đột giữa các mô hình quản lý và những rào cản nhận thức xã hội mang tính cố hữu. Việc thiếu vắng một công trình nghiên cứu hệ thống, khách quan và toàn diện dưới lăng kính khoa học giáo dục và thể chế học về toàn bộ chặng đường từ năm 1945 đến nay đã để lại những khoảng trống lý luận nhất định, khiến cho nhiều bài học kinh nghiệm xương máu của quá khứ chưa được tổng kết thấu đáo để phục vụ việc hoạch định chính sách hiện đại.

Trong bối cảnh đất nước bước vào kỷ nguyên phát triển mới dưới tác động mãnh liệt của cuộc Cách mạng công nghiệp lần thứ tư, sự bùng nổ của trí tuệ nhân tạo, yêu cầu chuyển đổi số và chuyển đổi xanh, hệ thống giáo dục nghề nghiệp Việt Nam đang đứng trước yêu cầu cấp thiết phải đổi mới căn bản và toàn diện theo tinh thần Nghị quyết số 29-NQ/TW của Ban Chấp hành Trung ương Đảng và Chỉ thị số 21-CT/TW của Ban Bí thư Trung ương Đảng. Do đó, việc thực hiện đề tài nghiên cứu \textbf{"Sự hình thành và phát triển hệ thống giáo dục nghề nghiệp Việt Nam từ năm 1945 đến nay"} không chỉ có giá trị lịch sử học thuần túy nhằm phục dựng bức tranh quá khứ, mà còn mang ý nghĩa thời sự đặc biệt sâu sắc trong việc đúc kết các quy luật vận động thể chế, nhận diện căn nguyên của các điểm nghẽn cố hữu, từ đó cung cấp luận cứ khoa học tin cậy cho công cuộc hoàn thiện hệ thống giáo dục nghề nghiệp Việt Nam trong giai đoạn 2026 -- 2035, tầm nhìn đến năm 2045.

\subsection*{2. Mục đích và nhiệm vụ nghiên cứu}

Mục đích tổng quát của báo cáo là nghiên cứu một cách có hệ thống tiến trình hình thành và phát triển của hệ thống giáo dục nghề nghiệp Việt Nam từ năm 1945 đến nay; phân tích sâu sắc các yếu tố tác động và đúc kết những bài học kinh nghiệm lịch sử có giá trị quy luật; từ đó đề xuất các định hướng và khuyến nghị chính sách khả thi nhằm hiện đại hóa hệ thống trong kỷ nguyên số.

Để đạt được mục đích tổng quát nêu trên, báo cáo tập trung giải quyết bốn nhiệm vụ nghiên cứu cụ thể:

Thứ nhất, hệ thống hóa và phân kỳ tiến trình lịch sử phát triển của hệ thống giáo dục nghề nghiệp Việt Nam từ năm 1945 đến nay qua sáu giai đoạn cụ thể, làm rõ sự biến đổi của bộ máy quản lý nhà nước, mạng lưới cơ sở đào tạo và cơ cấu trình độ đào tạo gắn liền với các văn bản quy phạm pháp luật nền tảng.

Thứ hai, phân tích đa chiều các nhóm yếu tố kinh tế, xã hội, chính trị tác động trực tiếp và gián tiếp đến sự phát triển của hệ thống qua từng thời kỳ lịch sử.

Thứ ba, đánh giá khách quan các thành tựu mang tính lịch sử, chỉ rõ những hạn chế cố hữu và nguyên nhân sâu xa; đúc kết bốn bài học kinh nghiệm có giá trị quy luật về tổ chức quản lý, cân đối cơ cấu nhân lực, gắn kết doanh nghiệp và phân luồng người học.

Thứ tư, nhận diện bối cảnh phát triển mới và đề xuất các nhóm giải pháp, khuyến nghị chính sách mang tính đột phá nhằm hoàn thiện hệ thống giáo dục nghề nghiệp Việt Nam đáp ứng yêu cầu công nghiệp hóa, hiện đại hóa và hội nhập quốc tế trong giai đoạn đến năm 2035.

\subsection*{3. Đối tượng và phạm vi nghiên cứu}

Xét về đối tượng nghiên cứu, đề tài tập trung khảo cứu hệ thống giáo dục nghề nghiệp Việt Nam trên các bình diện cơ bản: thể chế và bộ máy quản lý nhà nước, hệ thống văn bản quy phạm pháp luật, mạng lưới cơ sở đào tạo (các trường đào tạo công nhân kỹ thuật, trường dạy nghề, trung tâm dạy nghề, trường trung cấp và cao đẳng), cơ cấu trình độ đào tạo và mối quan hệ giữa đào tạo với thị trường lao động.

Về phạm vi thời gian, công trình giới hạn khảo sát từ tháng 9 năm 1945 (thời điểm thành lập nước Việt Nam Dân chủ Cộng hòa) đến năm 2026, với tầm nhìn dự báo chính sách đến năm 2035 và năm 2045. Về phạm vi không gian và nội dung, đề tài nghiên cứu trên phạm vi toàn lãnh thổ Việt Nam, tập trung trọng tâm vào giáo dục nghề nghiệp chính quy và các chương trình đào tạo nghề có tính chất định hướng chiến lược quốc gia.

\subsection*{4. Phương pháp luận và phương pháp nghiên cứu}

Đề tài được triển khai dựa trên cơ sở phương pháp luận của chủ nghĩa duy vật biện chứng và chủ nghĩa duy vật lịch sử, xem xét sự phát triển của giáo dục nghề nghiệp như một tiến trình vận động khách quan, luôn chịu sự chi phối và tương tác biện chứng với các điều kiện kinh tế, chính trị và xã hội của đất nước qua từng thời kỳ.

Trong quá trình thực hiện, báo cáo kết hợp đồng bộ các phương pháp nghiên cứu khoa học chuyên ngành:

Phương pháp nghiên cứu lịch sử kết hợp phương pháp logic được sử dụng để tái hiện tiến trình phát triển theo đúng trình tự thời gian khách quan của các sự kiện, đồng thời phát hiện ra bản chất, mối liên hệ bên trong và quy luật vận động của hệ thống qua các giai đoạn.

Phương pháp phân tích văn bản quy phạm pháp luật và văn kiện lịch sử được áp dụng để mổ xẻ các sắc lệnh, nghị định, quyết định của Chính phủ, luật của Quốc hội và nghị quyết của Đảng, làm rõ cơ sở pháp lý và định hướng đường lối của Nhà nước đối với đào tạo nghề.

Phương pháp so sánh lịch sử được vận dụng để xây dựng ma trận đối chiếu giữa các giai đoạn, làm nổi bật sự biến đổi về thể chế, mô hình tổ chức, cấp bậc đào tạo và cơ chế tài chính.

Phương pháp thống kê, tổng hợp và xử lý dữ liệu thứ cấp được sử dụng để khai thác các chuỗi số liệu thực chứng từ các báo cáo thường niên giáo dục nghề nghiệp quốc gia (Vietnam TVET Report) và báo cáo tổng kết của cơ quan quản lý nhà nước, đảm bảo tính xác thực tuyệt đối cho các nhận định khoa học.

\subsection*{5. Ý nghĩa khoa học và giá trị thực tiễn của báo cáo}

Về mặt khoa học, báo cáo góp phần hệ thống hóa cơ sở lý luận và thực tiễn về lịch sử phát triển giáo dục nghề nghiệp Việt Nam, bổ sung nguồn tư liệu học thuật có độ tin cậy cao phục vụ công tác giảng dạy, nghiên cứu trong các học phần thuộc ngành Khoa học giáo dục.

Về mặt thực tiễn, các phân tích phản biện và bài học kinh nghiệm lịch sử được đúc kết từ báo cáo cung cấp những luận cứ tham khảo có giá trị thiết thực cho các nhà hoạch định chính sách, cán bộ quản lý giáo dục nghề nghiệp trong quá trình sửa đổi, bổ sung Luật Giáo dục nghề nghiệp và triển khai Chiến lược phát triển giáo dục nghề nghiệp giai đoạn đến năm 2030, tầm nhìn đến năm 2045.
